% Luc Destombe --- licence CC BY-NC-SA 4.0
% https://creativecommons.org/licenses/by-nc-sa/4.0/deed.fr
% Fichier autonome : compiler avec pdflatex, deux fois (signets, nombre de pages).
\def\lyceeversion{1.9}
\documentclass[a4paper,10pt]{article}
\makeatletter
\newif\iflycee@niveaux
\newif\iflycee@palatino
\lycee@palatinotrue

\RequirePackage[utf8]{inputenc}
\RequirePackage[T1]{fontenc}
\RequirePackage[french]{babel}
\RequirePackage{etoolbox}

\newcommand{\ShorthandsOff}{\@ifundefined{shorthandoff}{}{\shorthandoff{;:!?}}}
\newcommand{\ShorthandsOn}{\@ifundefined{shorthandon}{}{\shorthandon{;:!?}}}
\ShorthandsOff
\AtBeginDocument{\ShorthandsOn}
\AtBeginEnvironment{tikzpicture}{\ShorthandsOff}

\RequirePackage{geometry}
\RequirePackage{fancyhdr}
\RequirePackage{lastpage}
\RequirePackage{multicol}
\RequirePackage{array}
\RequirePackage{enumitem}
\RequirePackage{xcolor}
\RequirePackage{needspace}

\RequirePackage{amsmath,amssymb}
\RequirePackage{mathpazo}
\renewcommand{\ttdefault}{lmtt}
\DeclareMathSizes{10.95}{10.95}{8}{5.5}
\begingroup\catcode`\_=\active \catcode`\^=\active
\gdef_#1{\sb{\scriptscriptstyle #1}}
\gdef^#1{\sp{\scriptscriptstyle\lycee@moinsepais #1}}
\endgroup
\newcommand\lycee@moins{\mathbin{%
  \setbox\z@\hbox{$\m@th\scriptscriptstyle\lycee@moinsorig$}%
  \dimen@=.55\wd\z@ \advance\dimen@-.5pt
  \hbox to.75\wd\z@{\hss$\m@th\scriptscriptstyle\vcenter{\hbox{%
    \tikz\draw[line width=.5pt,line cap=round](0,0)--(\the\dimen@,0);}}$\hss}}}
\begingroup\lccode`\~=`\- \lowercase{\endgroup\let~}\lycee@moins
\newcommand\lycee@moinsepais{\mathcode`\-="8000 }
\AtBeginDocument{\mathchardef\lycee@moinsorig=\mathcode`\-\relax}
\AtBeginDocument{\catcode`\_=12 \mathcode`\_="8000
  \catcode`\^=12 \mathcode`\^="8000 }
\newcommand\lycee@exposants{%
  \fontdimen13\textfont2=.48\fontdimen6\textfont2
  \fontdimen14\textfont2=.48\fontdimen6\textfont2
  \fontdimen15\textfont2=.48\fontdimen6\textfont2
  \fontdimen13\scriptfont2=.48\fontdimen6\scriptfont2
  \fontdimen14\scriptfont2=.48\fontdimen6\scriptfont2
  \fontdimen15\scriptfont2=.48\fontdimen6\scriptfont2}
\every@math@size\expandafter{\the\every@math@size\lycee@exposants}
\newlength{\retraitcalculs}
\setlength{\retraitcalculs}{2em}

\RequirePackage{tikz}
\usetikzlibrary{arrows.meta,calc,babel,shapes.misc}
\RequirePackage[most]{tcolorbox}
\tcbuselibrary{skins,breakable}

\definecolor{coulDef}{HTML}{1F4E79}
\definecolor{coulProp}{HTML}{7030A0}
\definecolor{coulRem}{HTML}{C55A11}
\definecolor{coulEx}{HTML}{2E7D32}
\definecolor{coulExo}{HTML}{B03A2E}
\definecolor{coulPart}{HTML}{1F4E79}
\definecolor{coulMeth}{HTML}{1B6E4F}
\definecolor{coulCourbe}{HTML}{0B5ED7}
\definecolor{coulCourbeB}{HTML}{C00000}
\definecolor{coulCourbeC}{HTML}{2E7D32}
\definecolor{coulGrille}{HTML}{8C8C8C}
\definecolor{coulRep}{HTML}{1B6E4F}
\definecolor{coulBar}{HTML}{2E7D32}

\newcommand{\R}{\mathbb{R}}
\newcommand{\Rs}{\mathbb{R}^{*}}
\renewcommand{\le}{\leqslant}
\renewcommand{\ge}{\geqslant}
\renewcommand{\approx}{\simeq}
\newcommand{\vect}[1]{\overrightarrow{#1}}
\newcommand{\norme}[1]{\left\lVert #1 \right\rVert}
\newcommand{\intervalle}[2]{\left[#1\,;#2\right]}
\RequirePackage{cancel}
\renewcommand{\CancelColor}{\color{coulExo}}
\newcommand{\fois}[1]{\times{\color{coulExo}#1}}
\newcommand{\barre}[1]{\cancel{#1}}

\tikzset{grille/.style={coulGrille,line width=0.4pt}}
\newcommand{\quadrillage}[4]{\draw[grille,step=1] (#1,#2) grid (#3,#4);}
\tikzset{
  croix/.style={cross out,draw,line width=0.6pt,inner sep=0pt,minimum size=4pt},
  croixdroite/.style={croix,rotate=45,line width=0.8pt},
}
\newcommand{\pt}[3]{\path (#1) node[croix]{} node[#2,font=\small,inner sep=2.5pt]{$#3$};}

\newcommand{\lycee@repere}[4]{%
  \draw[grille,step=1] (#1,#2) grid (#3,#4);
  \draw[-{Stealth[length=2mm]},thick] (#1,0)--({#3+0.4},0) node[below left=-1pt]{$x$};
  \draw[-{Stealth[length=2mm]},thick] (0,#2)--(0,{#4+0.4}) node[below left=-1pt]{$y$};
  \node[below left,font=\scriptsize,inner sep=1.5pt] at (0,0) {$0$};}
\newcommand{\lycee@gradx}[1]{\foreach \k/\l in {#1}{%
  \draw (\k,0.07)--(\k,-0.07) node[below,font=\scriptsize,inner sep=1.5pt]{$\l$};}}
\newcommand{\lycee@grady}[1]{\foreach \k/\l in {#1}{%
  \draw (0.07,\k)--(-0.07,\k) node[left,font=\scriptsize,inner sep=1.5pt]{$\l$};}}
\newcommand{\lycee@intff}[2]{\left[#1\,;#2\right]}
\newcommand{\lycee@intfo}[2]{\left[#1\,;#2\right[}
\newcommand{\lycee@intof}[2]{\left]#1\,;#2\right]}
\newcommand{\lycee@intoo}[2]{\left]#1\,;#2\right[}
\newcommand{\lycee@ens}[1]{\left\{#1\right\}}
\AtBeginDocument{%
  \providecommand{\repere}{\lycee@repere}%
  \providecommand{\gradx}{\lycee@gradx}%
  \providecommand{\grady}{\lycee@grady}%
  \providecommand{\Cf}{\mathcal{C}\sb{\scriptscriptstyle f}}%
  \providecommand{\Cg}{\mathcal{C}\sb{\scriptscriptstyle g}}%
  \providecommand{\intff}{\lycee@intff}%
  \providecommand{\intfo}{\lycee@intfo}%
  \providecommand{\intof}{\lycee@intof}%
  \providecommand{\intoo}{\lycee@intoo}%
  \providecommand{\ens}{\lycee@ens}}

\newlist{questions}{enumerate}{3}
\setlist[questions,1]{label=\arabic*\textdegree),leftmargin=*,itemsep=2pt,topsep=3pt}
\setlist[questions,2]{label=\alph*),leftmargin=*,itemsep=1pt,topsep=2pt}
\setlist[questions,3]{label=\roman*),leftmargin=*,itemsep=1pt,topsep=1pt}
\setlist[itemize]{label=\textbullet,leftmargin=*,itemsep=2pt,topsep=3pt}

\newcommand{\besoinplace}[1]{\par\penalty\@M\begingroup
  \dimen@=#1\relax\dimen@ii=\pagegoal\advance\dimen@ii-\pagetotal
  \ifdim\dimen@>\dimen@ii\ifdim\dimen@ii>\z@\vfil\fi\break\fi\endgroup}

\newcommand{\lycee@pied}{\thepage/\pageref{LastPage}}
\newcommand{\entete}[2]{%
  \pagestyle{fancy}\fancyhf{}%
  \fancyhead[L]{\footnotesize\color{coulPart}\textbf{#1}}%
  \fancyhead[R]{\footnotesize\color{coulPart}\textbf{#2}}%
  \fancyfoot[C]{\footnotesize\lycee@pied}%
  \renewcommand{\headrulewidth}{0.4pt}%
  \renewcommand{\headrule}{\hbox to\headwidth{%
    \color{coulPart!40}\leaders\hrule height \headrulewidth\hfill}}}

\RequirePackage{ccicons}
\newcommand{\lycee@licence}{{\scriptsize\color{gray}%
  \copyright\ Luc Destombe --- \ccbyncsa\ CC BY-NC-SA 4.0}}
\newif\iflycee@licence \lycee@licencetrue
\newcommand{\sanslicence}{\lycee@licencefalse}
\AtBeginDocument{\iflycee@licence\AddToHookNext{shipout/foreground}{%
  \put(0,0){\rlap{\kern\dimexpr1in+\hoffset+\oddsidemargin+\textwidth\relax
    \llap{\raisebox{-\dimexpr1in+\voffset+\topmargin+\headheight+\headsep
      +\textheight+\footskip\relax}{\lycee@licence}}}}}\fi}

\newcommand{\lycee@site}{faire-des-maths.fr}
\newcommand{\lycee@version}{{\scriptsize\color{gray}\lycee@site
  \ifdefined\lyceeversion\ --- Version~\lyceeversion\fi}}
\AtBeginDocument{\iflycee@licence\AddToHookNext{shipout/foreground}{%
  \put(0,0){\rlap{\kern\dimexpr1in+\hoffset+\oddsidemargin\relax
    \raisebox{-\dimexpr1in+\voffset+\topmargin+\headheight+\headsep
      +\textheight+\footskip\relax}{\lycee@version}}}}\fi}

\newcommand{\formulesserrees}{%
  \setlength{\abovedisplayskip}{3pt}\setlength{\belowdisplayskip}{3pt}%
  \setlength{\abovedisplayshortskip}{2pt}\setlength{\belowdisplayshortskip}{3pt}}

\setlength{\parindent}{0pt}

\RequirePackage[implicit=false,hidelinks,pdfencoding=auto,psdextra,bookmarksopen,
  bookmarksopenlevel=1,pdfcreator={},pdfproducer={}]{hyperref}
\RequirePackage{bookmark}
\hypersetup{pdfauthor={Luc Destombe},pdfkeywords={CC BY-NC-SA 4.0}}
\newcounter{lycee@signet}
\newcommand{\signet}[3][]{\stepcounter{lycee@signet}%
  \edef\lycee@dest{\if\relax\detokenize{#1}\relax
    signet.\arabic{lycee@signet}\else#1\fi}%
  \hypertarget{\lycee@dest}{}\bookmark[level=#2,dest=\lycee@dest]{#3}}
\pdfstringdefDisableCommands{%
  \def\R{R}\def\vect#1{#1}\def\Cf{Cf}\def\Cg{Cg}\def\fois#1{×#1}%
  \def\barre#1{#1}\def\intervalle#1#2{[#1;#2]}\def\intff#1#2{[#1;#2]}%
  \def\textsuperscript#1{#1}\def\dfrac#1#2{#1/#2}\def\frac#1#2{#1/#2}%
  \def\vec#1{#1⃗}}%
\RequirePackage[official]{eurosym}

\geometry{top=0.8cm,bottom=1.3cm,left=1cm,right=1cm,footskip=0.6cm}

\pagestyle{fancy}
\fancyhf{}
\renewcommand{\headrulewidth}{0pt}
\fancyfoot[C]{\footnotesize Page \thepage{} / \pageref{LastPage}}

\newcommand{\titrefiche}[2]{%
  \begin{center}
    \Large\textbf{#1}\\[2pt]
    \large\textbf{#2}\\[2pt]
  \end{center}
  \vspace{0.10cm}}

\setlength{\columnseprule}{0.4pt}
\setlength{\columnsep}{15pt}
\renewcommand{\columnseprulecolor}{\color{black!45}}
\setlength{\multicolsep}{2pt}

\newcounter{partiefiche}
\renewcommand{\thepartiefiche}{\Roman{partiefiche}}
\newif\ifapresbandeau
\newcommand{\lycee@bandeau}[1]{%
  \par\penalty-600
  \vspace{0.08cm}%
  \noindent\colorbox{coulPart}{%
    \parbox{\dimexpr\linewidth-2\fboxsep\relax}{%
      \color{white}\bfseries\raggedright\strut#1\strut}}%
  \nopagebreak\par\nopagebreak
  \vspace{0.06cm}\nopagebreak
  \apresbandeautrue}
\newcommand{\partiefiche}[1]{%
  \refstepcounter{partiefiche}%
  \lycee@bandeau{\signet{1}{\thepartiefiche{} --- #1}\thepartiefiche{} --- #1}}
\newcommand{\partiefichesansnum}[1]{\lycee@bandeau{\signet{1}{#1}#1}}

\newcounter{exo}
\newlength{\espaceexo}\setlength{\espaceexo}{7pt}
\newlength{\espacetitre}\setlength{\espacetitre}{2pt}
\newcommand{\exo}[1][\relax]{%
  \par
  \ifapresbandeau\apresbandeaufalse\else\penalty-150\addvspace{\espaceexo}\fi
  \refstepcounter{exo}%
  \noindent\signet[exercice-\arabic{exo}]{2}{Exercice \arabic{exo}}%
  \textbf{\color{coulExo}\underline{Exercice \theexo}}%
  \ifx\relax#1\relax\else\textbf{ : #1}\fi
  \par\nopagebreak\vspace{\espacetitre}\nopagebreak
  \ignorespaces}

\setlist[enumerate]{nosep,leftmargin=*,itemsep=2pt,topsep=2pt}
\setlist[itemize]{nosep,leftmargin=*,topsep=2pt}
\newcommand{\choix}[3]{\par\hspace*{1em}%
  \makebox[0.3\linewidth][l]{a) #1}\makebox[0.3\linewidth][l]{b) #2}c) #3}
\newcommand{\deux}[2]{\par\hspace*{0.6em}%
  \makebox[0.48\linewidth][l]{#1}#2}
\newcommand{\trois}[3]{\par\hspace*{0.6em}%
  \makebox[0.32\linewidth][l]{#1}\makebox[0.32\linewidth][l]{#2}#3}

\renewenvironment{center}{\par\addvspace{3pt}\centering}{\par\addvspace{3pt}}
\formulesserrees
\AtBeginDocument{\formulesserrees}
\clubpenalty=10000
\widowpenalty=10000
\displaywidowpenalty=10000
\brokenpenalty=10000
\setlength{\parskip}{0pt}

\RequirePackage{ragged2e}
\setlength{\RaggedRightRightskip}{0pt plus 1fil}
\AtBeginDocument{\RaggedRight\binoppenalty=10000 \relpenalty=10000 }
\g@addto@macro\@parboxrestore{\RaggedRight}
\makeatother

\tikzset{
  vecteur/.style={-{Stealth[length=2.4mm]},line width=1pt,coulCourbe},
  vecteurB/.style={-{Stealth[length=2.4mm]},line width=1pt,coulCourbeB},
  vecteurC/.style={-{Stealth[length=2.4mm]},line width=1pt,coulCourbeC},
  fig/.style={line width=0.8pt,black!75},
  etiq/.style={font=\small\bfseries,coulExo},
}

\begin{document}

\titrefiche{Seconde --- Vecteurs}{Fiche d'exercices du chapitre}

{\small\itshape Les figures se font au crayon et à la règle. On reproduit le
quadrillage quand il y en a un, et on vérifie chaque vecteur en comptant les
carreaux.\par}
\vspace{0.12cm}

\begin{multicols}{2}\raggedcolumns

\partiefiche{Caractéristiques d'un vecteur}

\exo[Égaux ou non]
La figure ci-dessous est tracée sur un quadrillage.
\begin{center}
\begin{tikzpicture}[scale=0.4]
  \quadrillage{-0.4}{-0.4}{11.4}{3.4}
  \draw[fig] (0,0)--(8,0) (3,3)--(11,3) (0,0)--(3,3) (4,0)--(7,3) (8,0)--(11,3);
  \pt{0,0}{below}{K} \pt{4,0}{below}{L} \pt{8,0}{below}{M}
  \pt{3,3}{above}{N} \pt{7,3}{above}{P} \pt{11,3}{above}{R}
\end{tikzpicture}
\end{center}
Tracer chaque vecteur sur la figure, puis compléter par $=$ ou $\neq$ :
\begin{enumerate}
  \item[a)] $\vect{KL}\;\ldots\;\vect{NP}$ \quad b)~$\vect{KN}\;\ldots\;\vect{LP}$
  \item[c)] $\vect{KP}\;\ldots\;\vect{LR}$ \quad d)~$\vect{NL}\;\ldots\;\vect{KN}$
  \item[e)] $\vect{KM}\;\ldots\;\vect{NR}$ \quad f)~$\vect{PR}\;\ldots\;\vect{PN}$
\end{enumerate}

\exo[Lire sur la figure]
On utilise la figure de l'exercice 1.
\begin{enumerate}
  \item Citer un vecteur égal à $\vect{KN}$, puis un vecteur égal à $\vect{LM}$.
  \item Citer un vecteur opposé à $\vect{KN}$, puis un vecteur opposé à $\vect{KL}$.
  \item Citer deux vecteurs :
  \begin{itemize}
    \item de même direction et de même sens, mais pas égaux ;
    \item de même norme, mais pas égaux ;
    \item de même direction et de sens contraires, mais pas opposés.
  \end{itemize}
\end{enumerate}

\exo[Construire un représentant]
Tracer un vecteur $\vec{u}$ de norme $3$~cm. Construire ensuite un représentant du
vecteur $\vec{v}$ dans chaque cas :
\begin{enumerate}
  \item $\vec{v}$ a la même direction et le même sens que $\vec{u}$, et pour norme
        $6$~cm ;
  \item $\vec{v}$ a la même direction et la même norme que $\vec{u}$, et le sens
        contraire ;
  \item $\vec{v}$ a la même direction que $\vec{u}$, le sens contraire, et pour norme
        $1{,}5$~cm ;
  \item $\vec{v}$ a la même norme que $\vec{u}$ et une autre direction.
\end{enumerate}

\exo[Dans un parallélogramme]
$RSTU$ est un parallélogramme de centre $O$. Faire une figure.
\begin{enumerate}
  \item Donner un vecteur égal à :
        \choix{$\vect{RS}$}{$\vect{TU}$}{$\vect{RU}$}
  \item Compléter par le point qui convient :
        \choix{$\vect{RO}=\vect{O\ldots}$}{$\vect{UO}=\vect{O\ldots}$}{$\vect{SO}=\vect{O\ldots}$}
\end{enumerate}

\partiefiche{Vecteurs et déplacements dans un quadrillage}

\exo[Lire des déplacements]
Pour chaque vecteur $\vec{a}$, $\vec{b}$, $\vec{c}$, $\vec{d}$ et $\vec{e}$, lire son
déplacement : tant de carreaux à droite ou à gauche, tant vers le haut ou vers le bas.
\begin{center}
\begin{tikzpicture}[scale=0.4]
  \quadrillage{-0.4}{-0.4}{13.4}{4.4}
  \draw[vecteur]  (0,0)--(3,4) node[midway,left,font=\small]{$\vec{a}$};
  \draw[vecteurB] (4,4)--(8,3) node[midway,above,font=\small]{$\vec{b}$};
  \draw[vecteurC] (8,0)--(4,1) node[midway,below,font=\small]{$\vec{c}$};
  \draw[vecteur]  (9,4)--(9,0) node[midway,left,font=\small]{$\vec{d}$};
  \draw[vecteurB] (13,0)--(10,3) node[midway,above right,font=\small]{$\vec{e}$};
\end{tikzpicture}
\end{center}

\exo[Neuf points]
Neuf points sont placés sur un quadrillage.
\begin{center}
\begin{tikzpicture}[scale=0.4]
  \quadrillage{-0.4}{-0.4}{6.4}{6.4}
  \pt{0,6}{above left}{A} \pt{3,6}{above}{B} \pt{6,6}{above right}{C}
  \pt{0,3}{left}{D} \pt{3,3}{above right}{E} \pt{6,3}{right}{F}
  \pt{0,0}{below left}{G} \pt{3,0}{below}{H} \pt{6,0}{below right}{I}
\end{tikzpicture}
\end{center}
\begin{enumerate}
  \item Citer tous les vecteurs formés par deux de ces points et égaux à $\vect{AB}$.
  \item Citer tous ceux qui sont égaux à $\vect{AE}$.
  \item Citer deux vecteurs opposés à $\vect{AB}$.
\end{enumerate}

\exo[Tracer des déplacements]
Reproduire les points, puis tracer :
\begin{enumerate}
  \item[a)] le vecteur $\vec{r}$ d'origine $A$ : $4$ carreaux à droite et $3$ vers le haut ;
  \item[b)] le vecteur $\vec{s}$ d'origine $B$ : $3$ carreaux à gauche et $2$ vers le haut ;
  \item[c)] le vecteur $\vec{t}$ d'origine $C$ : $4$ carreaux vers le bas ;
  \item[d)] le vecteur $\vec{m}$ d'origine $D$ : $2$ carreaux à droite et $4$ vers le bas.
\end{enumerate}
\begin{center}
\begin{tikzpicture}[scale=0.4]
  \quadrillage{-0.4}{-0.4}{13.4}{4.4}
  \pt{0,0}{below left}{A} \pt{8,1}{below right}{B}
  \pt{10,4}{above left}{C} \pt{11,4}{above right}{D}
\end{tikzpicture}
\end{center}

\exo[Origine donnée]
Reproduire chaque figure.
\begin{enumerate}
  \item À gauche : tracer un représentant de $\vec{u}$ d'origine $E$, puis un autre
        représentant de $\vec{u}$, d'origine $F$.
  \item À droite : tracer un représentant de $\vec{v}$ d'origine $G$, puis un autre
        représentant de $\vec{v}$, d'origine $H$.
\end{enumerate}
\begin{center}
\begin{tikzpicture}[scale=0.4]
  \quadrillage{-0.4}{-0.4}{7.4}{5.4}
  \draw[vecteur] (0,3)--(4,5) node[midway,above left,font=\small]{$\vec{u}$};
  \pt{0,0}{below left}{E} \pt{3,0}{below right}{F}
\end{tikzpicture}
\hspace{6pt}
\begin{tikzpicture}[scale=0.4]
  \quadrillage{-0.4}{-0.4}{7.4}{5.4}
  \draw[vecteurB] (0,5)--(3,2) node[midway,above right,font=\small]{$\vec{v}$};
  \pt{4,5}{above}{G} \pt{0,3}{left}{H}
\end{tikzpicture}
\end{center}

\exo[Extrémité donnée]
Reproduire chaque figure. Ici, chaque point nommé est la \textbf{pointe} de la flèche.
\begin{enumerate}
  \item À gauche : tracer un représentant de $\vec{u}$ d'extrémité $K$, puis un autre
        représentant de $\vec{u}$, d'extrémité $L$.
  \item À droite : tracer un représentant de $\vec{v}$ d'extrémité $M$, puis un autre
        représentant de $\vec{v}$, d'extrémité $N$.
\end{enumerate}
\begin{center}
\begin{tikzpicture}[scale=0.4]
  \quadrillage{-0.4}{-0.4}{6.4}{5.4}
  \draw[vecteur] (0,3)--(3,5) node[midway,above left,font=\small]{$\vec{u}$};
  \pt{4,2}{right}{K} \pt{6,5}{above}{L}
\end{tikzpicture}
\hspace{6pt}
\begin{tikzpicture}[scale=0.4]
  \quadrillage{-0.4}{-0.4}{6.4}{5.4}
  \draw[vecteurB] (0,5)--(1,2) node[midway,left,font=\small]{$\vec{v}$};
  \pt{4,1}{below}{M} \pt{6,2}{right}{N}
\end{tikzpicture}
\end{center}

\exo[Vecteurs opposés]
Reproduire la figure.
\begin{enumerate}
  \item Lire le déplacement de $-\vec{u}$, puis celui de $-\vec{v}$.
  \item Tracer un représentant de $-\vec{u}$ d'origine $P$, puis un représentant de
        $-\vec{v}$ d'origine $S$.
  \item Tracer un représentant de $-\vec{u}$ d'extrémité $R$.
\end{enumerate}
\begin{center}
\begin{tikzpicture}[scale=0.4]
  \quadrillage{-0.4}{-0.4}{11.4}{6.4}
  \draw[vecteur]  (0,5)--(4,6) node[midway,above left,font=\small]{$\vec{u}$};
  \draw[vecteurB] (10,6)--(11,3) node[midway,left,font=\small]{$\vec{v}$};
  \pt{5,2}{below right}{P} \pt{8,1}{below right}{S} \pt{4,4}{below left}{R}
\end{tikzpicture}
\end{center}

\partiefiche{Translation}

\exo[Image d'un quadrilatère]
Reproduire la figure, puis construire l'image $A'B'C'D'$ du quadrilatère $ABCD$ par la
translation de vecteur $\vec{u}$.
\begin{center}
\begin{tikzpicture}[scale=0.4]
  \quadrillage{-0.4}{-0.4}{9.4}{7.4}
  \draw[vecteur] (3,6)--(8,7) node[midway,above,font=\small]{$\vec{u}$};
  \draw[fig] (0,1)--(3,0)--(4,3)--(1,4)--cycle;
  \pt{0,1}{left}{A} \pt{3,0}{below}{B} \pt{4,3}{right}{C} \pt{1,4}{above left}{D}
\end{tikzpicture}
\end{center}

\exo[Dire avec les mots]
Reproduire la figure.
\begin{center}
\begin{tikzpicture}[scale=0.4]
  \quadrillage{-0.4}{-0.4}{12.4}{5.4}
  \draw[vecteur]  (0,2)--(3,5) node[midway,above left,font=\small]{$\vec{u}$};
  \draw[vecteurB] (5,5)--(9,4) node[midway,above,font=\small]{$\vec{v}$};
  \pt{1,0}{below}{A} \pt{5,2}{left}{C} \pt{7,0}{below}{E}
\end{tikzpicture}
\end{center}
\begin{enumerate}
  \item Placer le point $B$ tel que $\vect{AB}=\vec{u}$, puis compléter la phrase :
        « $B$ est l'image de \ldots{} par la translation de vecteur \ldots{} ».
  \item Placer le point $D$ tel que $\vect{CD}=\vec{v}$. Par quelle translation
        obtient-on $D$ à partir de $C$ ? Répondre par une phrase avec les mots
        « image », « translation » et « vecteur ».
  \item Placer le point $F$, image de $E$ par la translation de vecteur $\vec{u}$.
        Écrire l'égalité de vecteurs qui définit $F$.
  \item Par quelle translation obtient-on $C$ à partir de $D$ ? Répondre par une phrase.
\end{enumerate}

\exo[Quatre points à placer]
$ABC$ est un triangle. Placer les points $M$, $N$, $P$ et $R$ tels que
\[ \vect{AC}=\vect{BM}=\vect{NB}, \qquad \vect{PC}=\vect{AB}, \qquad \vect{AR}=\vect{CB}. \]

\exo[Origine ou extrémité]
On donne un vecteur non nul $\vec{u}$ et deux points $E$ et $F$. Tracer un
représentant de $\vec{u}$ d'origine $E$, puis un autre représentant d'\emph{extrémité}
$F$.

\exo[Image d'un triangle]
Tracer un triangle $ABC$ tel que $AB=5$~cm, $BC=4$~cm et $AC=3$~cm. Construire les
images $A'$, $B'$, $C'$ de ses sommets par la translation de vecteur $\vect{BC}$. Que
dire des triangles $ABC$ et $A'B'C'$ ?

\exo[Image ou point de départ]
On considère la translation de vecteur $\vect{CD}$.
\begin{center}
\begin{tikzpicture}[scale=0.4]
  \quadrillage{-0.4}{-0.4}{10.4}{6.4}
  \pt{0,4}{left}{C} \pt{3,1}{below}{D} \pt{4,6}{above}{A} \pt{5,3}{left}{B}
  \pt{10,2}{right}{E}
\end{tikzpicture}
\end{center}
\begin{enumerate}
  \item Reproduire la figure et tracer le vecteur $\vect{CD}$.
  \item Construire les images $A'$ de $A$ et $B'$ de $B$.
  \item Construire le point $F$ dont l'image est $E$. Écrire d'abord l'égalité de vecteurs
        qu'il vérifie.
\end{enumerate}

\partiefiche{Parallélogramme et égalité de vecteurs}

\exo[Construire, nommer, démontrer]
Reproduire chaque figure.
\begin{enumerate}
  \item À gauche : tracer le vecteur $\vect{AB}$, puis construire le point $D$ tel que
        $\vect{CD}=\vect{AB}$.
  \item À droite : tracer le vecteur $\vect{KL}$, puis construire le point $N$ tel que
        $\vect{MN}=\vect{KL}$.
\end{enumerate}
Dans chaque cas, tracer le parallélogramme obtenu, le nommer, puis démontrer que c'en
est un.
\begin{center}
\begin{tikzpicture}[scale=0.4]
  \quadrillage{-0.4}{-0.4}{5.4}{5.4}
  \pt{0,1}{left}{A} \pt{4,0}{below}{B} \pt{1,4}{above left}{C}
\end{tikzpicture}
\hspace{6pt}
\begin{tikzpicture}[scale=0.4]
  \quadrillage{-0.4}{-0.4}{6.4}{5.4}
  \pt{0,4}{left}{K} \pt{3,1}{below}{L} \pt{3,5}{above}{M}
\end{tikzpicture}
\end{center}

\exo[Deux parallélogrammes]
$EFGH$ et $EFIJ$ sont deux parallélogrammes : les points $E$ et $F$ sont communs aux
deux. Faire une figure.
\begin{enumerate}
  \item Donner, en justifiant, deux vecteurs égaux à $\vect{EF}$.
  \item Démontrer que $HGIJ$ est un parallélogramme.
\end{enumerate}

\exo[Autour d'un triangle]
Tracer un triangle $DEF$.
\begin{enumerate}
  \item Construire le point $M$ tel que $\vect{FM}=\vect{DE}$. Tracer le parallélogramme
        obtenu, le nommer, puis démontrer que c'en est un.
  \item Construire le point $N$ tel que $\vect{NF}=\vect{DE}$. Tracer le parallélogramme
        obtenu, le nommer, puis démontrer que c'en est un.
  \item Construire le point $K$, image de $D$ par la translation de vecteur $\vect{FE}$.
        Tracer le parallélogramme obtenu, le nommer, puis démontrer que c'en est un.
  \item Quel point de la figure a pour image $E$ par la translation de vecteur
        $\vect{FD}$ ? Justifier.
\end{enumerate}

\exo[Parallélogrammes accolés]
$EFGH$ et $HGKL$ sont deux parallélogrammes. Faire une figure, puis dire si chaque
affirmation est vraie ou fausse, en justifiant.
\begin{enumerate}
  \item L'image de $H$ par la translation de vecteur $\vect{GF}$ est le point $E$.
  \item Les vecteurs $\vect{KG}$ et $\vect{HL}$ ont la même direction.
  \item $\vect{KG}=\vect{HL}$.
  \item $\vect{EF}=\vect{LK}$.
\end{enumerate}

\exo[Vrai ou faux]
Dire si chaque affirmation est vraie ou fausse, en justifiant.
\begin{enumerate}
  \item Si $F$ est l'image de $C$ par la translation de vecteur $\vect{DE}$, alors
        $\vect{FC}=\vect{DE}$.
  \item Si $\vect{KL}=\vect{NM}$, alors $KLMN$ est un parallélogramme.
  \item Si $ABCD$ est un parallélogramme, la translation de vecteur $\vect{AD}$
        transforme $B$ en $C$.
\end{enumerate}

\exo[Quatrième sommet]
Reproduire la figure. Pour chaque point, écrire l'égalité de vecteurs que fournit le
théorème 1, puis construire le point.
\begin{center}
\begin{tikzpicture}[scale=0.4]
  \quadrillage{-0.4}{-0.4}{9.4}{8.4}
  \draw[fig] (2,4)--(6,3)--(5,7)--cycle;
  \pt{2,4}{left}{A} \pt{6,3}{below right}{B} \pt{5,7}{above}{C}
\end{tikzpicture}
\end{center}
\begin{enumerate}
  \item $D$ tel que $ABCD$ soit un parallélogramme.
  \item $E$ tel que $ABEC$ soit un parallélogramme.
  \item $F$ tel que $AFBC$ soit un parallélogramme.
\end{enumerate}

\partiefiche{Somme de deux vecteurs}

\exo[Origine imposée]
On pose $\vec{w}=\vec{u}+\vec{v}$. Reproduire chaque figure et tracer le représentant de
$\vec{w}$ d'origine $T$.
\begin{center}
\begin{tikzpicture}[scale=0.4]
  \quadrillage{-0.4}{-0.4}{7.4}{5.4}
  \pt{2,0}{below left}{T}
  \draw[vecteur]  (0,2)--(1,5) node[midway,left,font=\small]{$\vec{u}$};
  \draw[vecteurB] (3,5)--(6,4) node[midway,above right,font=\small]{$\vec{v}$};
\end{tikzpicture}
\hspace{6pt}
\begin{tikzpicture}[scale=0.4]
  \quadrillage{-0.4}{-0.4}{7.4}{5.4}
  \pt{4,5}{above}{T}
  \draw[vecteur]  (0,5)--(3,4) node[midway,below left,font=\small]{$\vec{u}$};
  \draw[vecteurB] (2,3)--(1,0) node[midway,left,font=\small]{$\vec{v}$};
\end{tikzpicture}
\end{center}

\exo[Trois points à construire]
Tracer un triangle $ABC$, puis construire les points $D$, $E$ et $F$ tels que
\[ \vect{BD}=\vect{BC}+\vect{CA}, \quad \vect{BE}=\vect{BA}+\vect{BC}, \quad
   \vect{CF}=\vect{CB}+\vect{AB}. \]

\exo[Sommes sur un quadrillage]
Reproduire chaque figure, puis placer le point $M$ tel que $\vect{AM}=\vec{u}+\vec{v}$.
\begin{center}
\begin{tikzpicture}[scale=0.4]
  \node[etiq] at (-1.2,5) {a)};
  \quadrillage{-0.4}{-0.4}{7.4}{5.4}
  \pt{1,3}{left}{A}
  \draw[vecteur]  (0,0)--(3,2) node[midway,above left,font=\small]{$\vec{u}$};
  \draw[vecteurB] (6,5)--(7,2) node[midway,left,font=\small]{$\vec{v}$};
\end{tikzpicture}
\hspace{6pt}
\begin{tikzpicture}[scale=0.4]
  \node[etiq] at (-1.2,5) {b)};
  \quadrillage{-0.4}{-0.4}{7.4}{5.4}
  \pt{0,1}{left}{A}
  \draw[vecteur]  (0,1)--(4,2) node[midway,above,font=\small]{$\vec{u}$};
  \draw[vecteurB] (5,0)--(7,3) node[midway,right,font=\small]{$\vec{v}$};
\end{tikzpicture}

\medskip
\begin{tikzpicture}[scale=0.4]
  \node[etiq] at (-1.2,5) {c)};
  \quadrillage{-0.4}{-0.4}{7.4}{5.4}
  \pt{0,1}{left}{A}
  \draw[vecteur]  (0,1)--(4,0) node[midway,below,font=\small]{$\vec{u}$};
  \draw[vecteurB] (0,1)--(2,4) node[midway,left,font=\small]{$\vec{v}$};
\end{tikzpicture}
\hspace{6pt}
\begin{tikzpicture}[scale=0.4]
  \node[etiq] at (-1.2,5) {d)};
  \quadrillage{-0.4}{-0.4}{7.4}{5.4}
  \pt{3,3}{below}{A}
  \draw[vecteur]  (0,5)--(3,3) node[midway,below left,font=\small]{$\vec{u}$};
  \draw[vecteurB] (5,0)--(6,3) node[midway,right,font=\small]{$\vec{v}$};
\end{tikzpicture}
\end{center}

\exo[Différences sur un quadrillage]
Reproduire chaque figure, puis placer le point $M$ tel que $\vect{BM}=\vec{u}-\vec{v}$.
\begin{center}
\begin{tikzpicture}[scale=0.4]
  \node[etiq] at (-1.2,5) {a)};
  \quadrillage{-0.4}{-0.4}{7.4}{5.4}
  \pt{1,4}{left}{B}
  \draw[vecteur]  (0,0)--(3,1) node[midway,above left,font=\small]{$\vec{u}$};
  \draw[vecteurB] (5,1)--(6,4) node[midway,right,font=\small]{$\vec{v}$};
\end{tikzpicture}
\hspace{6pt}
\begin{tikzpicture}[scale=0.4]
  \node[etiq] at (-1.2,5) {b)};
  \quadrillage{-0.4}{-0.4}{7.4}{5.4}
  \pt{0,2}{left}{B}
  \draw[vecteur]  (0,2)--(4,3) node[midway,above,font=\small]{$\vec{u}$};
  \draw[vecteurB] (5,1)--(7,4) node[midway,right,font=\small]{$\vec{v}$};
\end{tikzpicture}

\medskip
\begin{tikzpicture}[scale=0.4]
  \node[etiq] at (-1.2,5) {c)};
  \quadrillage{-0.4}{-0.4}{7.4}{5.4}
  \pt{1,2}{left}{B}
  \draw[vecteur]  (1,2)--(5,3) node[midway,below right,font=\small]{$\vec{u}$};
  \draw[vecteurB] (1,2)--(2,5) node[midway,left,font=\small]{$\vec{v}$};
\end{tikzpicture}
\hspace{6pt}
\begin{tikzpicture}[scale=0.4]
  \node[etiq] at (-1.2,5) {d)};
  \quadrillage{-0.4}{-0.4}{7.4}{5.4}
  \pt{3,3}{below}{B}
  \draw[vecteur]  (0,5)--(3,3) node[midway,below left,font=\small]{$\vec{u}$};
  \draw[vecteurB] (5,5)--(7,2) node[midway,left,font=\small]{$\vec{v}$};
\end{tikzpicture}
\end{center}

\exo[Avec des longueurs]
$ABC$ est un triangle tel que $AB=6$~cm, $AC=5$~cm et $BC=4$~cm. Faire une figure en
vraie grandeur.
\begin{enumerate}
  \item Construire le point $D$ tel que $\vect{AD}=\vect{AB}+\vect{AC}$.
  \item Construire le point $E$ tel que $\vect{BE}=\vect{BC}-\vect{BA}$.
  \item Quelle est la nature du quadrilatère $ABDC$ ? Justifier.
\end{enumerate}

\partiefiche{Multiplier un vecteur par un nombre réel}

\exo[Multiples d'un vecteur]
Reproduire la figure, puis représenter, à partir du point $O$, les vecteurs
$2\vec{a}$, $3\vec{a}$, $-\vec{a}$ et $-2\vec{a}$.
\begin{center}
\begin{tikzpicture}[scale=0.4]
  \quadrillage{-0.4}{-0.4}{11.4}{5.4}
  \pt{5,3}{above}{O}
  \draw[vecteur] (5,3)--(7,2) node[midway,above right,font=\small]{$\vec{a}$};
\end{tikzpicture}
\end{center}

\exo[Avec des coefficients]
$ABC$ est un triangle. Construire les points $M$, $N$ et $P$ tels que
\[ \vect{AM}=3\vect{AC}, \qquad \vect{CN}=\tfrac12\,\vect{CB}, \qquad
   \vect{BP}=-\tfrac12\,\vect{BA}. \]

\exo[Sur un quadrillage]
Quatre points $A$, $B$, $C$ et $D$ sont placés sur un quadrillage.
\begin{center}
\begin{tikzpicture}[scale=0.45]
  \quadrillage{-0.4}{-0.4}{8.4}{5.4}
  \pt{0,1}{below left}{A} \pt{4,1}{below}{B}
  \pt{6,3}{right}{C} \pt{1,4}{above left}{D}
\end{tikzpicture}
\end{center}
Reproduire la figure, puis tracer un représentant de chaque vecteur, d'origine
imposée (chaque extrémité tombe sur un n\oe ud du quadrillage) :
\begin{enumerate}
  \item[a)] $2\vect{BC}$, origine $A$ \quad b)~$0{,}5\,\vect{AC}$, origine $B$
  \item[c)] $-\tfrac13\,\vect{BD}$, origine $D$ \quad d)~$\tfrac13\,\vect{BD}$, origine $C$
  \item[e)] $\vect{AB}-\vect{BC}$, origine $D$ \quad f)~$\vect{AB}+2\vect{BC}$, origine $A$
\end{enumerate}

\exo[Multiples et sommes]
On pose $\vec{r}=2\vec{u}+\vec{v}$, $\vec{s}=\vec{u}-\vec{v}$ et
$\vec{t}=2\vec{u}-\vec{v}$. Reproduire la figure, puis tracer un représentant de
$\vec{r}$, de $\vec{s}$ et de $\vec{t}$ d'origine $O$.
\begin{center}
\begin{tikzpicture}[scale=0.4]
  \quadrillage{-0.4}{-0.4}{13.4}{6.4}
  \pt{0,1}{left}{O}
  \draw[vecteur]  (8,4)--(11,5) node[midway,above left,font=\small]{$\vec{u}$};
  \draw[vecteurB] (12,6)--(13,3) node[midway,left,font=\small]{$\vec{v}$};
\end{tikzpicture}
\end{center}

\partiefiche{Relation de Chasles}

\exo[Un seul vecteur, si c'est possible]
Écrire avec un seul vecteur lorsque c'est possible ; sinon, expliquer pourquoi.
\begin{enumerate}
  \item[a)] $\vect{BE}+\vect{EG}$ \quad b)~$\vect{DA}+\vect{DC}$ \quad c)~$\vect{TU}+\vect{KT}$
  \item[d)] $\vect{FH}+\vect{HS}$ \quad e)~$\vect{MN}+\vect{NM}$ \quad f)~$\vect{AB}+\vect{CE}$
  \item[g)] $\vect{AB}-\vect{CB}$ \quad h)~$\vect{KL}-\vect{KM}$ \quad i)~$-\vect{RS}+\vect{RT}$
  \item[j)] $\vect{AB}-\vect{BC}$
\end{enumerate}

\exo[Sans signe moins]
Écrire sans signe moins, en inversant les lettres.
\begin{enumerate}
  \item[a)] $-\vect{AB}$ \quad b)~$-\vect{MN}$ \quad c)~$-2\vect{CD}$
  \item[d)] $\vect{EF}-\vect{GH}$ \quad e)~$3\vect{KL}-\vect{JK}$
\end{enumerate}

\exo[Simplifier]
$ABC$ est un triangle ; $J$, $K$ et $M$ sont des points. Simplifier le plus possible à
l'aide de la relation de Chasles.\\
\emph{Aide :} ${-\vect{CB}=\vect{BC}}$ et ${\vect{AB}+\vect{AB}=2\vect{AB}}$.
\par\smallskip
{\renewcommand{\arraystretch}{1.5}%
\hspace*{1em}\begin{tabular}{@{}l@{\qquad}l@{}}
  $\vec{u}=\vect{BC}+\vect{CA}+\vect{AB}$ & $\vec{v}=\vect{CA}-\vect{CB}+\vect{AC}$\\
  $\vec{w}=\vect{MC}-\vect{MA}+\vect{CB}$ & $\vec{s}=\vect{CA}+\vect{AB}-\vect{BC}$\\
  $\vec{t}=\vect{MK}-\vect{MJ}+2\vect{KC}$ &
\end{tabular}\par}

\exo[Regrouper]
Simplifier en regroupant les vecteurs égaux.
\begin{enumerate}
  \item[a)] $\vect{AB}+\vect{AB}+\vect{AB}$ \quad b)~$\vect{EF}+2\vect{EF}$ \quad c)~$\vect{RS}-\vect{SR}$
  \item[d)] $\vect{DE}+\vect{EF}+\vect{DF}$ \quad e)~$\vect{GH}-\vect{KH}+\vect{GK}$
  \item[f)] $3\vect{AB}-\vect{BA}$
\end{enumerate}

\exo[Compléter]
Compléter par les points qui conviennent.
\begin{enumerate}
  \item $\vect{BC}=\vect{BA}+\vect{A\ldots}$
  \item $\vect{CB}+\vect{\ldots A}=\vect{CA}$
  \item $\vect{A\ldots}+\vect{CA}=\vect{CB}$
  \item $\vect{BC}+\vect{DA}+\vect{CD}=\vect{\ldots}$
\end{enumerate}

\besoinplace{14\baselineskip}
\exo[Des différences à simplifier]
Construire, à partir du point $A$, un représentant de chaque vecteur demandé.
\emph{On peut prolonger le quadrillage sur la feuille.}
\begin{center}
\begin{tikzpicture}[scale=0.4]
  \quadrillage{-0.4}{-0.4}{6.4}{3.4}
  \pt{0,1}{left}{A} \pt{4,0}{below}{B}
  \pt{5,3}{above right}{C} \pt{1,3}{above}{D}
\end{tikzpicture}
\end{center}
\begin{enumerate}
  \item[a)] $\vect{AC}-\vect{AB}$ \quad b)~$\vect{AD}-\vect{BD}$
  \item[c)] $\vect{AB}-\vect{AD}$ \quad d)~$\vect{AC}-\vect{BD}$
\end{enumerate}

\partiefiche{Démonstrations vectorielles}

\exo[Des égalités à démontrer]
Démontrer chaque égalité (pour tous points) en partant du membre de gauche.
\begin{enumerate}
  \item[a)] $\vect{EF}+\vect{GH}=\vect{EH}+\vect{GF}$ \quad
            b)~$\vect{RS}-\vect{RT}+\vect{UT}=\vect{US}$
  \item[c)] $\vect{OB}-\vect{OA}+\vect{BC}=\vect{AC}$
\end{enumerate}

\exo[Deux parallélogrammes à démontrer]
$ABC$ est un triangle.
\begin{enumerate}
  \item Construire le point $D$ tel que $\vect{BD}=\vect{BA}+\vect{BC}$.
  \item Démontrer que $\vect{AD}=\vect{BC}$ ; nature de $ADCB$ ?
  \item Construire le point $E$ tel que $\vect{AE}=\vect{AB}-\vect{AC}$.
  \item Démontrer que $\vect{BE}=\vect{CA}$ ; nature de $BEAC$ ?
\end{enumerate}

\exo[Exprimer avec deux vecteurs]
$ABC$ est un triangle.
\begin{enumerate}
  \item Exprimer $\vec{u}$ en fonction de $\vect{AB}$ et $\vect{AC}$ :
        \par\hspace*{1em}a)~$\vec{u}=\vect{CB}$
        \quad b)~$\vec{u}=3\vect{CB}+\vect{BA}$
  \item Exprimer $\vec{v}$ en fonction de $\vect{CA}$ et $\vect{BC}$ :
        \par\hspace*{1em}a)~$\vec{v}=\vect{BA}+\vect{CA}$
        \quad b)~$\vec{v}=\vect{AB}+2\vect{BC}+\vect{AC}$
\end{enumerate}

\partiefiche{Milieu d'un segment}

\exo[Un milieu à démontrer]
$KLMN$ est un parallélogramme.
\begin{enumerate}
  \item Faire une figure.
  \item Construire le point $P$ tel que $KLPM$ soit un parallélogramme.
  \item Démontrer que $M$ est le milieu de $[NP]$.
\end{enumerate}

\exo[Avec la relation de Chasles]
$KLM$ est un triangle.
\begin{enumerate}
  \item Construire le point $N$ tel que $\vect{KN}=\vect{KM}+\vect{LM}$.
  \item Démontrer que $M$ est le milieu de $[LN]$.
  \item Construire le point $P$ tel que $\vect{LP}=\vect{LK}+\vect{MK}$, puis démontrer
        que $K$ est le milieu de $[MP]$.
\end{enumerate}

\exo[Le double d'un vecteur]
\begin{minipage}[t]{0.6\linewidth}\raggedright
Reproduire la figure ci-contre.
\begin{enumerate}
  \item Construire le point $D$ tel que $\vect{AD}=2\vect{AB}$, puis démontrer que $B$
        est le milieu de $[AD]$.
  \item Construire le point $E$ tel que $\vect{CE}=2\vect{CB}$, puis démontrer que $B$
        est le milieu de $[CE]$.
  \item En déduire la nature du quadrilatère $ACDE$.
\end{enumerate}
\end{minipage}\hfill
\begin{minipage}[t]{0.37\linewidth}\vspace{0pt}\centering
\begin{tikzpicture}[scale=0.4]
  \quadrillage{-0.4}{-0.4}{6.4}{6.4}
  \pt{0,1}{left}{A} \pt{3,3}{right}{B} \pt{1,6}{above}{C}
\end{tikzpicture}
\end{minipage}

\exo[Trois points, deux milieux]
$ABC$ est un triangle. Les points $M$, $N$ et $P$ sont tels que
\[ \vect{AM}=\vect{AC}-\vect{AB}, \quad \vect{CN}=\vect{CB}-\vect{BA}, \quad
   \vect{BP}=\vect{AC}+\vect{BA}-\vect{CB}. \]
\begin{enumerate}
  \item Simplifier chaque égalité avec la relation de Chasles.
  \item Construire les points $M$, $N$ et $P$.
  \item Démontrer que $C$ est le milieu de $[AN]$, puis de $[BP]$.
  \item En déduire la nature du quadrilatère $ABNP$.
\end{enumerate}

\exo[Symétrie centrale]
$E$, $F$ et $G$ sont trois points. $E'$ est le symétrique de $E$ par rapport à $G$, et
$F'$ celui de $F$ par rapport à $G$.
\begin{enumerate}
  \item Démontrer que $FEF'E'$ est un parallélogramme.
  \item Justifier que $\vect{E'F'}=\vect{FE}$, puis comparer les longueurs $E'F'$ et $EF$.
  \item Quelle propriété de la symétrie centrale vient-on de démontrer ?
\end{enumerate}

\end{multicols}

\end{document}
